% =====================================================================
%  Hans Brøchner, „Det ubevidste og det bevidste menneskelige Sjæleliv“.
%  Nyt dansk Maanedsskrift, udg. af Vilhelm Møller, 3. Bind (Oktober
%  1871 -- Marts 1872), printed pp. 42--57.
%  TRANSCRIPTION (Danish) -- diplomatic, from the page images.
%
%  ---- SOURCE -----------------------------------------------------------
%  Internet Archive item nyt-dansk-maanedsskrift-bind-3 (whole volume,
%  Nyt_dansk_Maanedsskrift_bind3.pdf, 595 PDF pp., sha256
%  b79c303a...e1d01a6d, ABBYY FineReader text layer), downloaded 2026-09-27:
%  https://archive.org/details/nyt-dansk-maanedsskrift-bind-3
%  The article is PDF pp. 48--63 of it.  Working scan: an extract of those
%  16 pages (qpdf, text layer kept),
%  ~/bibliotek/Brøchner, Hans/1871-ubevidste-bevidste-nyt-dansk-maanedsskrift-3.pdf
%  sha256 63090c423840ba28b29decde30e1707c581d12dbbe3720a61dd3798094900a1d.
%  Each page one ~305 ppi RGB JPEG.  Fraktur.
%
%  ---- PAGE MAP ---------------------------------------------------------
%  Volume PDF = printed + 6; extract PDF = printed - 41 (p. 42 = extract 1).
%  Verified by eye on every folio, pp. 42--57, 2026-09-27.
%
%  ---- CONVENTIONS ------------------------------------------------------
%  Diplomatic; 1871 orthography as printed.  Letterspacing -> \emph;
%  antiqua against the Fraktur -> \textit.  Fraktur I/J resolved by the
%  word (Individet, Idee; Jeg).  Quotation marks „...“ as printed.
%  Footnotes *) at the mark.  Printer's errors transcribed as printed and
%  logged in a % comment at the site.
%
%  ---- LAYOUT -----------------------------------------------------------
%  Head in bold Fraktur (p. 42); sections numbered 1. (p. 42), 2. (p. 46),
%  3. (p. 52), 4. (p. 55), centred; an unnumbered short rule before the
%  closing paragraph (p. 57).  Signed „H. Bröchner.“ -- with ö, as printed
%  (confirmed at 400 dpi) -- set right in bold Fraktur, then a double rule.
%  No footnotes.  No antiqua or bold in running text; 56 letterspaced
%  runs, measured within the line (spaced letter gaps 6--14 px at 300 dpi
%  against 1--5 px set solid).
%  Printer's errors logged at the site: p. 53 „Bevisthedsfordobling“;
%  p. 57 a gap about one short word wide before „Døden“ (a sort lost?).
%
%  ---- ACCURACY  [2026-09-27] ------------------------------------------
%  Word-accuracy: both batches diffed against the embedded ABBYY layer
%  before splicing (ocrdiff.py --frag, --offset -41): 2 batches (pp.
%  42-49, 50-57), 19 candidates, 0 transcription errors corrected, 19 the
%  witness's fault, 0 unresolved.  See TRANSCRIPTION-PLAYBOOK.md §3 step 4.
%  The transcribers took their words from that layer corrected at the
%  image, so the diff cannot see a layer error they copied; the second
%  reading exists for that class.
%  SECOND READING of all 16 pages, word by word and mark by mark at the
%  image (300 dpi, 600 dpi for any doubt), by two readers who had
%  transcribed nothing and were forbidden the text layer: 1 discrepancy.
%  p. 57 „eet Værende“ -> „eet Bærende“: the capital is a Fraktur B
%  (reader's template match against same-page capitals, B 0.93-0.94 vs
%  V 0.84-0.87; confirmed at 600 dpi in the caller; cf. „baaret“ below).
%  An inherited layer error -- the layer reads „Værende“ -- and the one
%  class the diff cannot see.  Everything else (words, punctuation,
%  emphasis, paragraphing, heads, close) agreed.  Result: 1 word error in
%  16 pages found and corrected.
% =====================================================================
\documentclass[12pt,a4paper,openany]{book}
\usepackage[utf8]{inputenc}
\usepackage[T1]{fontenc}
\usepackage{textalpha}
\usepackage{libertinus}
\usepackage{libertinust1math}
\usepackage{graphicx}    % for \reflectbox (the old Danish ɔ: id-est mark)
\DeclareUnicodeCharacter{0254}{\reflectbox{c}}  % ɔ (U+0254) = reversed c
\usepackage[danish]{babel}
\usepackage[protrusion=true,expansion=true]{microtype}
\usepackage{setspace}
\usepackage[margin=1.2in]{geometry}
\usepackage{fancyhdr}
\setlength{\headheight}{28pt}
\addtolength{\topmargin}{-13.5pt}

% Footnotes as the 1877 printing sets them: *).
\renewcommand{\thefootnote}{*)}

% The article's numbered sections: the number alone, centred on its line.
\newcommand{\secnum}[1]{\par\begin{center}#1\end{center}\par\nobreak}

\usepackage{hyperref}
\hypersetup{
  pdftitle={Det ubevidste og det bevidste menneskelige Sjæleliv},
  pdfauthor={Hans Brøchner},
  hidelinks
}

\pagestyle{fancy}
\fancyhf{}
\fancyhead[LE,RO]{\thepage}
\fancyhead[RE]{\textit{H. Brøchner: Det ubevidste og det bevidste menneskelige Sjæleliv}}
\fancyhead[LO]{\textit{Nyt dansk Maanedsskrift, 3.\ Bind}}

\setstretch{1.3}

\title{\textbf{Det ubevidste og det bevidste\\ menneskelige Sjæleliv}}
\author{Hans Brøchner}
\date{Nyt dansk Maanedsskrift, udg.\ af Vilhelm Møller,
  3.\ Bind (Oktober 1871--Marts 1872), pp.~42--57}

\usepackage{marginnote}
\newcommand{\opage}[1]{\marginnote{\footnotesize\textit{[#1]}}}
\newcommand{\apage}[1]{\marginnote{\footnotesize\textit{[#1]}}}

\begin{document}
\maketitle
\thispagestyle{empty}
\clearpage

% --- p. 42 ---
\opage{42}\begin{center}\textbf{\large Det ubevidste og det bevidste menneskelige Sjæleliv.}\end{center}
% p. 42: above the head stand the last lines of „Engelsk Familieliv“ (faintly printed) -- not transcribed. A short rule separates the head from „1.“.

\secnum{1.}

Blandt alle de Problemer, med hvilke Psykologien i vore
Dage med fornyet Interesse beskæftiger sig, er der intet, der
har et videre Omfang og en større Rækkevidde end Problemet
om Forholdet mellem vort bevidste og vort ubevidste Sjæleliv.
Den Række af Fænomener, det berører, naaer ligefra Individets
første Tilbliven til dets Ophør, og de Livsformer,
hvis Forhold det fordrer opklaret, vise tilbage til selve Tilværelsens
Grundmodsætninger og til Forholdet mellem disse.
Hvad der bestandig med uimodstaaelig Magt fører Erkjendelsen
tilbage til dette Problem, er vor Aands uafviselige
Trang til at finde Eenhed i os selv og i Tilværelsen, til
ikke at blive staaende ved Deeltheden som det Sidste. Det
er denne Trang, der i et Aarhundrede to Gange og i heelt
modsatte Retninger har ført til eensidige Formler for Menneskets
og Tilværelsens Eenhed: til Materialismens Eenhedsformel,
der lader det Aandelige i Tilværelsen forsvinde i den
bevægede Materie og det bevidste Liv i os blive et Resultat
af Stof-Elementernes Forbindelse; --- og til Spiritualismens
Eenhedsformel, der lader Naturen forsvinde i den naturløse
Aand og Tankelivet i os løsrive sig fra sin organiske Betingethed.
Det er denne Trang, der atter, efterat de eensidige
Formlers Uigjennemførlighed er bleven indlysende, fører til
% --- p. 43 ---
\opage{43}en Stræben efter Eenhed, som ikke miskjender Modsætningernes
Eiendommelighed, men vil see denne fremgaae af og atter
sammensluttes i et høiere, rigere Forenende. ---

Det var ved den nyere Filosofis Begyndelse, at den
skarpe, dualistiske Sondring mellem det, som vor Tids Psykologi
% emphasis: „Sjæleliv“ letterspaced (letter gaps 7--11 px at 300 dpi; plain words on the line 2--5 px)
betegner som det ubevidste og det bevidste \emph{Sjæleliv}
og stræber at føre tilbage til Eenhed, blev fastslaaet gjennem
en Konsekvens, der droges, ikke af en omfattende, uhildet
% „Erfaring,“: a faint speck follows the comma (the layer reads «); not a second point.
Erfaring, men af metafysiske Forudsætninger. Aand og
Materie, Tanke og Udstrækning, vare bestemte som Modsætninger,
% UNSURE: „Prædikater;“ -- the upper dot of the semicolon is faint and touches the lower stroke; could be a comma. The layer reads ;.
der ikke mødtes i fælles Prædikater; men kun
udelukkede hinanden. I Mennesket havde man dog, tvungen
af Kjendsgjerningernes Magt, maattet sætte Aand og Materie,
Sjæl og Legeme, som forbundne og som paavirkende hinanden;
men Forbindelsen bestemtes udtrykkelig som ubegribelig og
kun som en Sammensætning, ikke som en virkelig Gaaen
sammen til Eenhed. Alt, hvad der hører Sjælen til, skulde
gaae ind under Tankens og Bevidsthedens Bestemmelse, Alt,
hvad der hører Legemet til, under Udstrækningens og Bevægelsens,
fra hvilken Tanke og Bevidsthed udelukkedes. Legemet
blev derfor, om end Sjælen skulde have et „Sæde“ deri ---
men som en fremmed Gjæst og ved et ubegribeligt Under ---
kun opfattet som en Maskine, som en „guddommelig Automat“,
og Alt, hvad der foregaaer deri, skulde skee efter Mekanismens
Love. I Dyret, hvor der ikke anerkjendtes nogen Sjæl,
fordi Dyret frakjendtes Tanke og Bevidsthed, saae man dette
Forhold træde tydeligst frem. Den Vexelvirkning, der, om
end kun gjennem en Inkonsekvens, sattes i Mennesket mellem
Sjæl og Legeme, og som dog kun kunde virkeliggjøres i
Mekanismens Former, faldt her bort, og enhver Tilskyndelse
til en Virksomhed udgik her fra et Legemligt. Haardnakket
vægrede man sig da ved at opfatte selv de Fænomener indenfor
Dyreverdenen, der allermeest erindrede om den menneskelige
Bevidstheds, som Yttringer af en sjælelig Virksomhed. Naar
saaledes det mishandlede Dyr ved sit Skrig udtrykte sin
Smerte, saa benægtede man, i Tillid til sine Forudsætnin%
% --- p. 44 ---
\opage{44}gers Urokkelighed, dristigt, at Skriget var et Tegn paa en
Fornemmelse. Dyrets Skrig, sagde man, var som Strengens
Tone, en reen mekanisk Virkning af Slaget, der ramte det.

Fra denne dualistiske Theori, der sætter en uoverstigelig
Grændse mellem Menneskelivet og de lavere Livsformer, og
som i Menneskevæsenet vil forene det absolut Modsatte, altsaa
Uforenelige, til Eenhed og lade det til absolut forskjellige
Sfærer Henførte --- der altsaa maatte mangle enhver Mulighed
til en Berøring --- træde i Vexelvirkningsforhold, maatte dog
snart Erkjendelsen af de absurde Konsekvenser og af de
uløselige Vanskeligheder, til hvilke den ledte, føre bort, og et
friere, mere uhildet Blik paa Virkelighedens Fænomener og
en skarpere Prøvelse af Modsætningernes Begreber maatte
føre til at udvide Theoriens Bestemmelser af Aanden og
Naturen og derigjennem bane Veien for Erkjendelsen af den
Eenhed, der sammenslutter dem i Tilværelsens Heelhed og
i Menneskelivets Former.

Det var fra Aandens Bevidsthed om sig selv som tænkende,
at hiin dualistiske Theori havde taget sit Udgangspunkt, da
den søgte et sikkert Fodfæste for Erkjendelsen, og idet Aanden,
for at finde sin egen Væsensbestemmelse, sondrede sig fra Alt,
hvorfra den kunde abstrahere, blev det Legemlige udskilt
fra den som heelt fremmed for den og som det, der ingen
Bestemmelse havde tilfælles med den. Derfor blev Tanke,
og hvad der gjordes til Et dermed: Bevidsthed, ved en
uoverstigelig Kløft holdt ude fra Aandens Modsætning; det
sjælelige Liv fik kun een Form: Bevidsthedslivet, og mellem
dette Bevidsthedsliv og det, der foregik i det Ubevidste, kunde
ingen Eenhed tænkes. Og dog viser Bevidstheden med Nødvendighed
hen til en Væsenssammenhæng med det Ubevidste.
I dens Begreb ligger der et Væsens-Forhold til det, der er
et Andet end Jeget, i hvilket Andet ogsaa det Materielle er
omfattet; men et væsenligt Forhold indeslutter, som et nødvendigt
Forhold, Væsenssammenhængen; kun ved den forklares
Forholdets Nødvendighed. Og Bevidsthedslivet har
en Vorden, og idet det vorder, fremgaaer det af en ubevidst
% --- p. 45 ---
\opage{45}Leven, der er blivende Betingelse for Bevidsthedens Virkelighed.
Ved selve det Begreb, der skulde holde det Aandelige
ude fra det Legemlige, vises der saaledes hen til en Eenhed
i Menneskevæsenet af Aand og Natur, af det Bevidste og
det Ubevidste; og der vises derved hen til en fyldigere og
rigere Bestemmelse af Modsætningerne end den, der kunde
gives dem som abstrakt adskilte. For at Eenheden skal kunne
virkeliggjøres for Tanken, maa Naturens Væsen ikke bestemmes
ved det, der kun udtrykker det Reale (det Legemlige) ligesom
i dets største Afstand fra Aanden, ved den passive Udstrækning;
dens Begreb maa ikke dannes ud fra dens fattigste Former,
men ud fra de høieste Fænomener, i hvilke den rigest og klarest
udtrykker sit fuldstændige Væsen: fra det Organiske, det Levende,
i hvilket Naturens Idealitet kommer tilsyne. Og Aandens
Begreb maa ikke bestemmes alene ud fra de Fænomener, hvor
Sammenhængen med Naturgrundlaget ligesom er bleven usynlig,
men ud fra hele Kredsen af dens Former og derfor
ogsaa ud fra dem, der vise det Aandelige i dets Tidsudvikling
som baaret og betinget ved et realt organisk Grundlag,
det bevidste Aandsliv i Eenhed med et ubevidst sjæleligt, gjennem
hvilket dets Sammenhæng med Naturen bliver aabenbar.
Der vil da, fra begge Udgangspunkter, fra Naturen og fra
Aanden, naaes til den væsenlige Eenhed af Modsætningerne,
og kun ved Anerkjendelsen af denne væsenlige Eenhed ville
Kjendsgjerningerne kunne forstaaes efter deres Fuldstændighed.

Naar Modsætningerne i Menneskevæsenet sees i Eenhed,
% emphasis: „Livets“ letterspaced (letter gaps 7--11 px; plain words on the line 2--5 px)
saa bliver \emph{Livets} Begreb det, der forbinder dem; det er
paa een Gang det, der udtrykker Naturens Væsen og det,
der danner Grundlaget for Aandens. Vi see da Livet som
omfattende de to Former: det organiske og det i snevrere
Betydning sjælelige, det ubevidste og det bevidste Liv, og idet
Menneskevæsenet sees som Eenhed og Eenheden bestemmes fra
den høieste Forms Synspunkt, saa bliver Bestemmelsen af
% emphasis: „sjæleligt“, „dets Heelhed“, „Sjælen“, „al“ letterspaced (gaps 6--11 px); „Liv“, „Livet“, „Livsvirk-“ set solid (2--4 px)
\emph{sjæleligt} Liv i videre Forstand udstrakt til Livet i \emph{dets
Heelhed}, og \emph{Sjælen} sees som Princip for \emph{al} Livsvirksomhed
i det Menneskelige. Sjælen er da Et med Livs%
% --- p. 46 ---
\opage{46}principet, og dette er atter det ideelle Eenhedsprincip, ved
hvilket Organismen er som Totalitet, som en Eenhed, der
hævder sig selv i Forholdet til det Ydre og til sine Organer
og Funktioner i deres Mangfoldighed. Vi opfatte altsaa
ikke Sjælen som et Fremmed, der ligesom er en Beboer af
Legemet, eller som et Immaterielt, der paa en ubegribelig
Maade behersker Stof-Elementerne i Legemet; men vi opfatte
% emphasis: „egen“ letterspaced (gaps 7--10 px)
den som den reale Organismes \emph{egen} Idealitet, og dens
Idealitet sætte vi ikke som eensbetydende med en absolut
Modsætning til det Materielle, men som Formen for Eenhedens
% emphasis: „som“ letterspaced (gaps 9, 7 px; the line's other words 1--5 px)
Væren \emph{som} Eenhed i det reale Mangfoldige, der
netop ved at kunne blive Led i en Totalitet, i en levende
% emphasis: „som saadan“ and „enkelte“ letterspaced (gaps 7--14 px)
Heelhed, der er tilstede \emph{som saadan} i sine \emph{enkelte} Led,
viser sig at bære Idealitetens Mulighed i sig.

Men naar vi i Begrebet: sjæleligt Liv indeslutte begge
Livets Grundformer som i en høiere Eenhed, saa ligger der
heri, at vort bevidste og vort ubevidste Liv maa staae i et
gjennemgaaende Vexelforhold til hinanden i hele Rækken af
det individuelle Sjælelivs Fænomener, og dette gjennemgaaende
Vexelforhold, gjennem hvilket Væsenseenheden overalt hævder
sig, skulle vi nu med korte Træk søge at paavise, idet vi
forlade de abstrakte Begrebsbestemmelsers Sfære for at gaae
over til Kjendsgjerningernes.

\secnum{2.}

Som den første Virksomhed af det Livsprincip, der, som
% emphasis: „sjæleligt“ and „organiserende“ letterspaced (gaps 7--12 px)
Princip for begge Livsformer, bestemtes som \emph{sjæleligt} Princip,
finde vi den \emph{organiserende} Virksomhed. Som frembringende
en organisk Totalitet, ikke en Atomernes Sammenhobning,
viser denne Virksomhed sig som en ideel, sjælelig
Virksomhed, men som en saadan, der heelt er i det Ubevidstes
Form. Den er Forudsætningen for det bevidste
Sjæleliv, idet den uddanner de Organer, ved hvilke Bevidsthedslivets
Virkelighed er betinget. Men den er ikke blot
Tidsforudsætningen for Bevidsthedslivet; den er ogsaa, efterat
dette er kommet til Virkelighed, den blivende Betingelse
% --- p. 47 ---
\opage{47}for dette, idet den bestandig gjenfrembringer Bevidsthedslivets
Organer.

Til denne organdannende, reent ubevidste sjælelige Virksomhed
slutter sig saa med umærkelige, kontinueerlige Overgange,
Bevidsthedslivets elementæreste Former i Fosteret og
i det selvstændige Individ paa dets første Livstrin. Vi tage
her Bevidsthedslivets Begreb i dets videste Betydning, i
hvilken det ikke blot omfatter det Stadium, paa hvilket Jeget
klart sondrer sig som Jeg fra hvad der ikke er Jeget, men
ogsaa det Stadium, paa hvilket Bevidstheden endnu kun
dæmrer under Fornemmelsen. Fra de ubevidste Reflexbevægelser
gjennem de dunkleste Former af Organfornemmelser og Sandsefornemmelser,
med deres mekaniske Fremtrængen og Tilbagetrængen,
og til de bevidste Fornemmelser, og atter fra disse
til de mere og mere udfoldede Bevidsthedsfunktioner af Forestilling
og Tænkning, have vi en Række, i hvilken Overgangen
skeer gjennem det uendeligt Lille, og hvor de kvalitative
Forskjelles Grændse ikke lader sig drage ad Erfaringens Vei.
Paa hvert Punkt i den fremskridende Række er der et relativt
Ubevidst og et relativt Bevidst: den ubevidste Begyndelse
stræber hen imod det Bevidste, og dette gjemmer i sig det
Ubevidste som Rest eller som Element.

Hvorledes det Bevidsthedsliv, der paa denne Maade
udvikler sig af det Ubevidste, bevarer sin uopløselige Sammenhæng
med dette, kan fra mange forskjellige Synspunkter
gjøres tydeligt. Vi see saaledes, at der ikke blot tilfældigt,
% emphasis: „en Vexlen“ letterspaced (gaps 7--10 px; „en“ at the start of the next line 2 px)
men normalt og periodisk finder \emph{en Vexlen} Sted mellem
en Tilstand, i hvilken vor Bevidsthed er vaagen, og en Tilstand,
i hvilken vort sjælelige Liv er sænket ned i det Ubevidste.
I den drømfrie Søvn, Søvnens normale Form, finder denne
vort totale Bevidsthedslivs periodiske Tilbagegang i det Ubevidste
Sted paa en saadan Maade, at denne Tilbagegang viser sig
som ubetinget nødvendig, som overvældende enhver Villiens
Modstand. Og selve Bevidsthedslivets Sundhed, dets Samling
og dets Klarhed, er betinget ved denne periodiske Neddukken
i det Ubevidstes styrkende Bad.

% --- p. 48 ---
\opage{48}Vi see det samme Vexelforhold i det Fænomen,
% emphasis: „at ethvert ... af dette“ letterspaced, from „at“ at the end of the first line (a--t 8 px; „at“ set solid elsewhere on the page 2 px) through „dette“ (gaps 8--12 px)
\emph{at ethvert
bestemt Bevidsthedsindhold gaaer tilbage i
det Ubevidste og atter fremgaaer af dette}. Hvad
der i det enkelte Øieblik udfylder Bevidstheden, kan ikke blivende
være dens Indhold; ethvert Forsøg paa, ud over en vis Tidsgrændse
at fastholde det Enkelte i Bevidstheden, vil uundgaaelig
mislykkes; det trænges ligesom med Magt bort derfra.
Og at det maa det, ligger i Sagens Natur: det Enkelte er
netop som Enkelt inkongruent med Bevidsthedens alt det
Enkelte omfattende Almindelighed; Bevidstheden kan derfor
ikke presses ind i dets snevre Grændser. Men det, der
bliver draget bort fra Bevidsthedens Brændpunkt og snart
rykkes heelt ud af dens Synskreds, gaaer ikke tabt for
Bevidstheden. Det gaaer tilbage i det ubevidste Dyb og
gjemmes der som Aandens Eiendom. Og ikke blot gjemmes
det, saaledes at det under bestemte Betingelser atter kan
hæves op til Lyset og erindres, men det drages ind i en
Proces, der foregaaer i Selvet under Bevidsthedens Niveau:
det indgaaer i Forbindelser med det øvrige Bevarede, det
smeltes sammen med det og nuanceres derved. Der arbeides
uafbrudt i det Skjulte, og en Udvikling kan finde Sted,
der overrasker, idet dens Resultater komme frem for Bevidstheden,
som naar under Søvnen en Tanke er modnet, en
Beslutning klaret, eller naar under en længere aandelig
Hviletid en aandelig Assimilationsproces er foregaaet, det
spredte og fremmede Stof samlet og tilegnet. Denne Virksomhed
i Sjælens ubevidste Dyb sender bestandig fine Virkninger
op selv til de høieste og frieste Tankens og Villiens
Virksomheder, og den bevidste Aand har ingen Magt til
at afbryde det Arbeide og den Udvikling, der uophørlig
foregaaer i den tilsyneladende bevægelsesløse ubevidste Inderlighed.

Fra et nyt Synspunkt see vi den væsenlige Sammenhæng
% emphasis: „Bevægeligheden af Grændsen imellem dem“ letterspaced (gaps 6--11 px); „iagttage“ set solid
mellem de to Livsformer, idet vi iagttage \emph{Bevægeligheden
af Grændsen imellem dem}. Saaledes bliver
f. Ex. en Bevægelse, en Handling, der, naar den er uvant for
% --- p. 49 ---
\opage{49}Individet, kun udføres, idet Opmærksomheden fæstes paa de
bestemte Organdele, som medvirke til dens Udførelse, gjennem
Vanen udført ubevidst, paa mekanisk Maade (f. Ex. Behandlingen
af et Instrument). Kun det almindelige Billede af
Resultatet, og det endda kun bestandig mere fordunklet,
staaer for Bevidstheden: alle de Trin, gjennem hvilke Resultatet
tilveiebringes, blive os ubevidste. Indenfor visse Grændser
gjælder dette Sidste om enhver vilkaarlig, og som saadan bevidst
Bevægelse. Bevægelsens endelige Form staaer for Bevidstheden
som et Bevægelsesbillede; men Mellemleddene, ved
hvilke Bevægelsen realiseres: de Nervevirkninger og Muskelsammentrækninger,
der ere dens Betingelser, ere os ubevidste.
Omvendt kan det, der i vor Organismes normale Tilstand
er unddraget Bevidstheden, under abnorme Forhold blive os
bevidst: Grændsen altsaa bevæges i modsat Retning.

% emphasis: „den enkelte Bevidsthedsakt“ letterspaced (gaps 7--11 px; „vi“ before it 4 px)
Analysere vi \emph{den enkelte Bevidsthedsakt}, finde vi
atter det Bevidste og det Ubevidste gaaende i Eenhed, idet
selve Bevidsthedsakten er ubevidste Elementers Eenhed. Vælge
vi til Exempel den enkelte bevidste Sandsning, saa viser dette
sig let. Den enkelte Tonefornemmelse for Exempel, der for
vor umiddelbare Opfattelse staaer som et reent Enkelt, har
til sine Forudsætninger en fysisk og en organisk Proces, i
hvilken Elementerne ere lette at paavise. I det fysiske Incitament,
der betinger den bestemte Tonefornemmelse, er der et
Antal af Bølger, svarende til Tonens Eiendommelighed, og
i hver af disse Bølger, under deres Fremadskriden, utallige
Grader i Fortætning. Og i den organiske Proces, der fremkaldes
ved hiint Incitament, maa der være et Tilsvarende til
Incitamentets Elementer og Gradforskjelle. Men idet den
% emphasis: only „fornemmelse“ (after the line break) is letterspaced (gaps 7--12 px); „Sandse-“ at the end of the line above is set solid (3--5 px). Marked as printed.
organiske Proces omsættes i den psykiske, i den bevidste Sandse\emph{fornemmelse},
sammenfattes det Mangfoldige til en Eenhed,
i hvilken Elementerne ikke blive bevidste som saadanne, men
kun gjøre sig gjældende gjennem Fornemmelsens Kvalitets-Bestemthed.
Hvad der gjælder om den enkelte Sands, gjælder
om alle, forsaavidt Bevægelsesprocesser ere de forskjellige
Fornemmelsers fysiske og organiske Betingelser. I høiere
% --- p. 50 ---
\opage{50}Former af Bevidsthedslivet træffe vi atter det samme Forhold.
Idet vi danne en Almeenforestilling f. Ex. af et Træ, af et
Dyr, saa danne vi den uvilkaarlig, ved en mekanisk Proces,
i hvilken de enkelte Sandsningsbilleder blive Elementerne,
som frembringe et udpræget Billede der, hvor de ligesom dække
hinanden, men udviskes, hvor de ere forskjellige. Men denne
Proces, der er Forudsætningen for den bevidste Udhæven af
Begreberne, bliver tilligemed sine Elementer i det Ubevidste;
kun Resultatet kommer frem for Bevidstheden. Det er dette
Forhold mellem den enkelte Bevidsthedsakt og dens ubevidste
% emphasis: „Leibnitz“ letterspaced (gaps 9--11 px; „først“ before it 3--5 px)
Elementer, paa hvilket først \emph{Leibnitz} har henledet Opmærksomheden.
I Havets Brusen, hvilken vi opfatte som en samlet
Lyd, viste han hen paa de utallige enkelte Bølgers Lyd, der
uden at blive os bevidste samle sig til Eenhed for vort Øre;
i vor Opfattelse af den grønne Farve, der opstaaer ved
Blandingen af fiint revet Blaat og Guult, viste han hen
paa de uendelig smaa ubevidste Virkninger af de enkelte
Farve-Elementer i Blandingen. Overhovedet viste han hen til
det uendeligt Lille som det ubevidste Element i enhver bestemt
bevidst Størrelse.

Selve dette Forhold mellem Elementet og Elementernes
summerede Eenhed viser atter hen til et nyt Synspunkt for
Forholdet mellem det Ubevidste og det Bevidste. Hvad der
ifølge sin Kvantitets-Bestemthed bliver ubevidst for os, kan
paa en dobbelt Maade hæves til Bevidsthed. Det fysiske Incitament,
der fremkalder en Bevægelse i Organismen --- hvilken
vi ikke kunne tænke os adskilt fra en psykisk Akt: et Fornemmelseselement
--- men ikke en Bevægelse, der er stærk nok
% emphasis: „bevidst“ letterspaced (gaps 10--12 px; „en“ before it 4 px); „en Tilvæxt i Styrke“ letterspaced from „en“ at the end of the line (e--n 11 px) through „Styrke“ (gaps 7--11 px); „udløse“ solid
til at fremkalde en \emph{bevidst} Fornemmelse, kan ved \emph{en
Tilvæxt i Styrke} udløse en saadan. Det Samme kan
indenfor bestemte Grændser naaes ad en anden Vei: gjennem
% emphasis: „Opmærksomheden“ letterspaced (gaps 7--13 px)
en Henvendelse af \emph{Opmærksomheden}. Vi koncentrere
os bevidst i den bestemte Retning, og det Indtryk, der
før var for svagt til at blive ændset, et flygtigt Lydindtryk
f. Ex., opfattes nu klart og bevidst.

% --- p. 51 ---
% emphasis: „bestandig gjensidig Indvirkning“ letterspaced from „be-“ at the end of the line (b--e 12 px) through „Indvirkning“ (gaps 7--13 px); „Sted“ solid (3--5 px)
\opage{51}Vi see endelig Sammenhængen mellem det ubevidste og
det bevidste Sjæleliv fra det Synspunkt, at der finder en \emph{bestandig
gjensidig Indvirkning} Sted, gjennem hvilken
den væsenlige Eenhed aabenbarer sig. Og herigjennem see
% emphasis: „Kredsløb“ letterspaced (gaps 8--12 px)
vi et \emph{Kredsløb} tilveiebragt. Hvad der fra det bevidste
Livs Sfære ligesom nedfældes i det ubevidstes og bestemmer
dette, tildeels med synlige Virkninger i det Organiske, virker
derfra atter tilbage paa Bevidsthedslivet gjennem den usondrede,
koncentrerede sjælelig-aandelige Eenheds ubevidste Virken.
Ud fra den Bevidsthedslivets Resultater optagende, koncentrerede
Eenhed bestemmes ubevidst Aandsvirksomhedernes almindelige
Retning, og der udgaaer fra den, paa uvilkaarlig
Maade, Impulser til eiendommelige Forestillingssammenknytninger
(Ideeassociationer) og Supplementer og Korrektiver for
Erkjendelsen og Villien. Og fra hiint Ubevidste, til hvilket
Bevidsthedslivets Momenter, som trædende ud af Bevidsthedssfæren,
give deres Bidrag, og i hvilket de sammenarbeides,
% emphasis: „Tilbagevirkning“ letterspaced, including „Til-“ at the end of the line (gaps 9--11 px); „som Bevidsthedslivets Forudsætning og Grundlag“ letterspaced from „som“ at the end of the line (gaps 9--12 px; „allerede“ before it 3--5 px) through „Grundlag“ (gaps 8--13 px); „virker“ solid (2--4 px)
finder der ikke blot under disse Former en uophørlig \emph{Tilbagevirkning}
paa det Bevidste Sted; men allerede \emph{som
Bevidsthedslivets Forudsætning og Grundlag} virker
fra først af det ubevidste Liv ind paa det bevidste gjennem
de i dets individuelle Eiendommelighed begrundede Stemninger
og Følelser, Attraaer og Drift-Tilskyndelser, og paa særegen
Maade gjennem den Sammenhæng, i hvilken det bringer det
bevidste Liv med det universelle Livs Strømninger, med
hvilke det staaer i nærmere og inderligere Forbindelse. Af
den fortsatte Assimilation gjennem det Ubevidstes Proces afhænger
det bevidste Livs Udvikling; af denne Processes normale
Bevægelse Bevidsthedslivets Sundhed og uforstyrrede Harmoni.

Saaledes vises vi overalt hen til Eenheden af de to
Livsformer, til deres indre Væsensforbindelse. Paa ethvert
Udviklingsstadium vise de hen til hinanden: det Ubevidste til
det Bevidste som sit Maal, sin Fuldendelse og sin Eenhed;
det Bevidste til det Ubevidste som sin Forudsætning, sin blivende
Betingelse, sit Element. Forskjellige i Omfang --- idet Bevidsthedslivet
hviler paa det ubevidste Livs bredere Basis, saaledes
% --- p. 52 ---
% emphasis: „hele“ letterspaced (gaps 7--10 px)
\opage{52}at vort Jeg, skjøndt det betegner vor \emph{hele} Væren efter dens
reale og ideale Side, dog kun i den relative Abstraktions
Form udtrykker vor totale individuelle Leven --- ere de Et i
Væsen, Former for det samme sjælelige Princips Virksomhed.

\secnum{3.}

Idet vi i det Foregaaende betragtede de to Livsformers
Forhold, have vi seet dette alene fra det individuelle Livs
Synspunkt og som visende tilbage til Forholdet mellem
Natursiden og Aandssiden i os. Men Forholdet vil vise
sig i et nyt og klarere Lys, naar Synspunktet udvides fra
Individet til det i Individet virksomme Universelle, og naar
vi see Vexelforholdet mellem det Bevidste og det Ubevidste som
% emphasis: „Aands“ letterspaced (gaps 8--10 px); in „universellere aandelig“ only „aandelig“ is letterspaced (gaps 6--10 px), „universellere“ solid (2--6 px)
betinget ved den individuelle \emph{Aands} Forhold til en universellere
\emph{aandelig} Magt.

Vi naae allerede til dette Forhold ved at gaae tilbage
til Begrebet om Selvets koncentrerede Eenhed, der paa umiddelbar
Maade gjør sig gjældende i Individets bevidste Liv.
I denne koncentrerede Eenhed kunne vi nemlig skjelne mellem
% emphasis: „Erhvervede“ (gaps 7--10 px) and „Oprindelige“, spaced on both sides of the line break (gaps 7--13 px), letterspaced
det \emph{Erhvervede}, der, som assimileret, i den umiddelbare
Værens Form er gaaet ind i Eenheden, og det \emph{Oprindelige},
der er givet som Væsensanlæg og fra Begyndelsen af
virker ubevidst; og dette oprindelig ubevidst Virkende søge
% emphasis: „høieste aandelige“ (gaps 8--11 px), „Aprioriskes“ (gaps 6--11 px) and „Geniale“ (gaps 7--13 px) letterspaced
vi nu netop i de \emph{høieste aandelige} Funktioner: i det
\emph{Aprioriskes} Virken og i det \emph{Geniale}, idet vi, i hiint
som i dette, gjennem det Individualiteten inderligst Tilhørende
see et Universelt at gjøre sig gjældende.

Hvad vi, med et i Filosofien vedtaget Navn, benævne
% emphasis: „Aprioriske“ letterspaced (gaps 5--10 px)
som det \emph{Aprioriske}, betegner de oprindelige Former, i
hvilke vort aandelige Væsens Eiendommelighed giver sig et
Udtryk under al dets Virksomhed i Aandslivets forskjellige
Retninger, og saaledes giver sig et Udtryk, at disse Former
vise sig som forud for Paavirkningen anlagte i det Væsen,
der ved Paavirkningen bestemmes til Virksomhed. Paa ethvert
Gebeet i Menneskeaanden, hvor en eiendommelig Bevidsthedsvirksomhed
fremtræder, maa et eiendommeligt udpræget Aprio%
% --- p. 53 ---
\opage{53}risk gjøre sig gjældende: i vor Anskuen, i vor Tænken, i
vor ethiske Handlen og vor kunstneriske Frembringen. Men
dette Aprioriske, der udtrykker vort Væsens Grundeiendommelighed,
viser, idet dette Væsen ikke er ved sig selv, men er
begrundet i et Høiere, tilbage til dette: Tilværelsens Grundkilde,
og det bliver --- da det, der er vort Væsens Grund, maa
være væsensligt med det Begrundede, --- Udtryk (om end abstrakt
Udtryk) for selve Tilværelsesprincipets Væsen. Dette
Aprioriske, der er Væsensbestemmelse i os, bliver --- netop som
% emphasis: „Aanden“ (gaps 8--12 px), „selv“ at the end of the line (gaps 8--11 px), „erhverve“ (gaps 7--11 px) and „givet“ (gaps 10--13 px) letterspaced; „skal“ between „selv“ and „erhverve“ is solid (3--5 px), so marked as two runs
Væsensbestemmelse i \emph{Aanden}, der ifølge sin Frihed \emph{selv}
skal \emph{erhverve} hvad der oprindelig er den \emph{givet} --- kun ved en
Aandsvirksomhed til for Bevidstheden: det maa blive til for
den ved vor egen, af Erfaringen kun inciterede Energi. Men
inden det kommer til Bevidsthed, virker det normerende og
bestemmende i enhver Bevidsthedsakt; det forbereder ved sin
% printer's error: „Bevisthedsfordobling“ (d missing) as printed
skjulte Virken Muligheden af den Bevisthedsfordobling, ved
hvilken det selv kan blive Gjenstand for Bevidstheden, og i
selve den Akt, ved hvilken det træder frem for Bevidstheden,
% emphasis: „enhver“ letterspaced (gaps 10--12 px)
er det selv, som normerende \emph{enhver} Aandsakt, ubevidst bestemmende.
Det er vort Væsens Lov, der bestemmer den
Akt, ved hvilken det virksomme Væsens Grundformer først
skulle hæves til Bevidsthed. Men disse Væsenets Former
ere tillige et Høieres Former, Udtryk for det universelle
Princips Væsen, og derved gyldige ikke blot for det individuelle
Selv, men for Menneskeheden og for Tilværelsen.

Hvad der gjælder om det Aprioriske overhovedet, gjælder
med det Samme om den prægnante Form af det, der træder
% emphasis: „Geniale“ letterspaced (gaps 7--10 px)
frem i det \emph{Geniale}. Der er i dette en oprindelig Mulighed
i det Ubevidstes Form, der uimodstaaelig arbeider hen
imod en Aabenbaring i Bevidstheden; men selv naar Geniet
bliver sig klarest bevidst, bevares et Ubevidst dog i den
geniale Virken. Den Eenhed af Nødvendighed og Frihed,
af Naturvirken og Aandsenergi, der karakteriserer Geniet,
medens det instinktmæssig søger Aabenbaring, karakteriserer
det endnu i dets høieste Bevidsthedsudvikling. Dets høieste
Virken er i Umiddelbarhedens Form, en ubevidst Skaben,
% --- p. 54 ---
\opage{54}der først naar den har givet sit Værk Virkelighed, gjennem
det udstraaler det Lys, hvori den selv bliver deelviis gjennemsigtig.
Og i den geniale Frembringen --- i hvilken Individet
paa den intensiveste Maade udpræger sin aandelige Eiendommelighed
--- føler Individet sig som Organ for en høiere
Magt, af hvilken det gjennemvirkes. Geniets Gjerning har
Ideens Gyldighed; men derved viser den tilbage til den
ideelle Magt, fra hvilken Geniets Gave har sit Udspring.

Ved det Oprindeligste og Dybeste i den individuelle
% emphasis: „det Universelles Virken i Individet“ letterspaced from „det“ (d--e 10 px, e--t 8 px) through „Individet“ (gaps 6--12 px); „til“ before it solid
Aands Virken føres vi saaledes tilbage til \emph{det Universelles
Virken i Individet}, og her træder saa Forholdet mellem
det Bevidste og det Ubevidste i os frem i en eiendommelig
Belysning. Det Universelles Virken bliver, om end en
aandelig Virken, i Forhold til Individet et partielt eller totalt
Ubevidst, baade fordi det Universelle er det Omfattende,
det Overgribende, der ikke rummes i den individuelle Aand,
hvis Forhold til den er det enkelte Totalitetsleds Forhold til
selve Totaliteten og dens Princip, --- og fordi det Universelles
Virksomhedsformer ikke kongruere med den endelige
Bevidstheds Virksomhedsformer. Selv naar den individuelle
Aand klarest bliver sig sit Forhold til det gjennemvirkende
Universelle bevidst, bliver den ubevidste Virken kun indenfor
en vis Grændse omsat i Bevidsthed, fordi selve denne
Akt ubevidst styres og bestemmes af det dybere Virkende,
ligesom den Akt, ved hvilken det Aprioriske gjordes til Gjenstand
for Bevidsthed, bestemtes ved selve det Aprioriskes, Individet
ubevidste Virken.

Mellemleddet, gjennem hvilket Tilværelsesprincipets ubevidste
Indvirken paa det individuelle Bevidsthedsliv virkeliggjøres,
% emphasis: „er den reale Tilværelsesheelhed“ letterspaced, „er“ included (e--r 8 px; „er“ solid at the end of the same line 2 px), „gjøres“ solid
\emph{er den reale Tilværelsesheelhed}. Denne er,
som principbestemt Totalitet, i sin reale Væren det Fornuftbestemte,
% emphasis: „ubevidst“ letterspaced three times (gaps 8--14 px); the „det“ before each is solid
Planmæssige (Teleologiske); men den er det \emph{ubevidst}
Fornuftige, det \emph{ubevidst} Tænkende, det \emph{ubevidst} Hensigtsmæssige.
Den ubevidste Naturfornuft er Grundlaget for det
bevidste, individuelle Fornuftliv. Vi tænke kun, idet vi have
en Hjerne, der ved den ubevidste Naturtanke er formet til
% --- p. 55 ---
\opage{55}vor Tankes Organ. Bag vor bevidste Tanke ligger „det,
% p. 55: a stray mark stands before „og“ after „os“;“ -- not a letter, not transcribed
der tænker i os“; og den af Væsenet, af „Tankegeniet“
fremgaaende, uvilkaarlige, ubevidste Tankeproduktion, det, vi
kunde kalde „den indre Tænken“, „Tanken i og for sig“,
ligger bagved Bevidstheden om denne Tænken, bagved Fremstillingen
af Tanken, ikke blot for Andre, men for os selv.

% emphasis: „den historiske Udvikling“ letterspaced (gaps 7--15 px); „Under“ and „træder“ solid
Under \emph{den historiske Udvikling} træder Vexelforholdet
mellem det Bevidste og det Ubevidste frem i prægnante
Former. Idet den historiske Udvikling er en teleologisk Udvikling,
men i en saadan Udvikling paa ethvert Trin Udviklingens
% emphasis: „ideelt“ letterspaced (gaps 9--13 px)
Maal \emph{ideelt} er tilstede i sin Heelhed som det, der
bestemmer og leder Udviklingen: saa er Udviklingens styrende
og ledende Magt tilstede som det, der naaer ud over den
enkelte Bevidsthed, og hvis Indhold kun i en begrændset Form
er realiseret for Bevidstheden. Efter sin Heelhed er den derfor
et ubevidst Virkende: den virker bagved Bevidstheden,
skjult for Bevidstheden. Derfor virker den som Individets
% emphasis: „Skjæbne“ (gaps 8--12 px), „enkelt“ twice (gaps 8--12 px) and „abstrakt“ (gaps 6--10 px) letterspaced
\emph{Skjæbne}, der vel kan gjennemlyses af Bevidstheden, men
kun partielt, idet Tankens Princip bliver Gjenstand for
Tanken efter en \emph{enkelt} Bestemthed, et \emph{enkelt} Grundforhold,
eller kun paa \emph{abstrakt} Maade. Den historiske Udviklings
Princip, i sin Virken som konkret Totalitet, bevarer, som det
Individets Bevidsthed Gjennemvirkende, sit Præg af en
Skjæbne, der kun relativ forklares til bevidst Aandslov.
Det universelle Aandsprincip sammenslutter i hvert Moment,
ubevidst for Individet, den individuelle Bevidsthed under sin
overgribende Nødvendighed.

\secnum{4.}

Endnu fra eet Synspunkt viser det bevidste Sjælelivs
% emphasis: „Døden“ letterspaced (gaps 8--11 px)
Forhold til det ubevidste sig for os, idet vi see \emph{Døden} som
det totale Bevidsthedslivs Tilbagegang i det Ubevidste. Idet
% emphasis: „er“ at the end of the line letterspaced (e--r 8 px; „er“ solid elsewhere 2--3 px); „organisk“ at the start of the next line solid
vort Bevidsthedsliv i alle dets Former gjennemgaaende \emph{er}
organisk betinget, maa Organismens Opløsning nødvendigviis
medføre alle disse Bevidsthedsformers Ophør, og det
nærmeste Udtryk herfor bliver: en Tilbagegaan i det Ube%
% --- p. 56 ---
\opage{56}vidste. Men dette Udtryk kan skjule en dobbelt Betydning.
Det kan betegne en Opløsning af Selvet i de materielle
Elementer, hvis Kræfter samvirkede i det levende Individ.
Det vilde være Materialismens Fortolkning af Udtrykket.
Men det kan ogsaa betegne en Tilbagegaan i en ligesom
mere fordybet og mere inderliggjort Værens Form, for hvilken
Bevidsthedslivets Bestemmelser ikke mere have Gyldighed. Og
til en saadan Fortolkning vises vi tilbage derved, at Menneskeaandens
Frihed indeslutter Garantien for dens Evighed,
sikkrer den imod at kunne forsvinde som et Væsenløst; men
at den samme Frihed viser tilbage til det Guddommelige
som sin Kilde, til Selvhengivelsen til Principet, til Væsensforeningen
med det som det høieste Maal, og netop derved
viser ud over den Sfære, hvor Bevidsthedslivet med dets
Relativitet og endelige Modsætninger faaer en Plads. Men
naar vi saa spørge om Formen for det aandelige Livs Væren
efter at Bevidsthedslivet er forsvundet med sine Betingelser, saa
opkaste vi et Spørgsmaal, der, netop fordi vi i vor Tænken
ere bundne til Bevidsthedslivets Former, kun kan besvares
af os gjennem Antydninger og Analogier. Vi søge da en
Analogi i et Forhold, der i det Foregaaende blev berørt: i
den enkelte Bevidsthedsakts Forhold til det totale ubevidste
Sjæleliv og til Bevidsthedslivets Heelhed. Naar det enkelte
Bevidsthedsindhold, der er fremkaldt ved en ydre Paavirkning,
et Sandsningsbillede f. Ex., træder bort fra Bevidsthedens
Synskreds og gaaer tilbage i det Ubevidste, saa træder det
gjennem dette ind i en omfattende Sammenhæng af Erindringer
og Tanker; det bliver tilegnet af denne, omformet, inderliggjort
fra Sandsebillede til Forestilling og Tanke, og tilhører
som saadan vor omfattende aandelige Livseenhed. I denne
Inderliggjørelse gjennem Tilbagegangen i det Ubevidste vilde
vi kunne søge en Analogi, der kan antyde Formen for det
individuelle Selvs Forvandling ved Bevidsthedslivets Ophør.
Dets Gaaen tilbage i det Ubevidste vilde da være Et med
en Omformning fra den fænomenale, sandseligt bestemte Udvortesheds
Form til en dybere Inderlighed. Det vilde da
% --- p. 57 ---
\opage{57}være et Moment i det Eviges Liv, i den Eenhed, af hvilken
Livsudfoldelsen i Tilværelsen evigt fremgaaer; og det vilde
sammensluttes med de andre individuelle Aander i den evige
Livsfyldes Organisation. Ogsaa til Forholdet mellem de
individuelle Aander lader hiin Analogi sig udstrække. Ligesom
vore enkelte Bevidsthedsakter i deres bevidste Sondring
% emphasis: „under Bevidsthedens Grændselinie“ letterspaced (gaps 7--12 px); „Kontinuitet“ and „og“ solid
fra hverandre have en Kontinuitet \emph{under Bevidsthedens
Grændselinie} og derfor samles i Aandens omfattende,
sammensluttende Eenhed, saaledes kunne de enkelte individuelle
Aander, hvis Organismer jo ere Led i een omfattende fysisk
Totalitet, forestilles som sammensluttede i deres ubevidste
Grund af eet Bærende, Omfattende og Forbindende, og som
% printer's defect: the line „Døden gaaende tilbage ...“ begins with a blank about one short word wide (a sort apparently lost at the line start, e.g. „i“); transcribed as printed
Døden gaaende tilbage til denne Grund, og i den til en
dybere, inderligere Forening.

% p. 57: a short centred rule (c. 2 cm) stands before the closing paragraph; no section number
\begin{center}\rule{2cm}{0.4pt}\end{center}

Saaledes vilde da Bevidsthedslivet blive som et Omraade
i vort sjælelige Liv, der til begge Sider var begrændset
af det Ubevidste. Det hæver sig op af den ubevidste
% emphasis: „organiske“ (gaps 8--10 px) and „aandelig“ (gaps 8--10 px) letterspaced
\emph{organiske} Levens Grund, er, under det menneskelige Livs
Udvikling, i hvert Øieblik baaret og betinget af det Ubevidste,
i bestandig Vexelvirkning med det, og idet det naaer sin
Afslutning, gaaer det tilbage i et Ubevidst, der nu udtrykker
en høiere \emph{aandelig} Livets Form, hævet over Bevidsthedslivets
Begrændsning, over dets Udvorteshedssondring og dets
Endelighed: et Liv i Eenhed med det Guddommelige, der i
sin evige Viden af sig ikke som det bevidste Selv begrændses
af et Fremmed, men i denne Viden indeslutter Tilværelsens
Heelhed, og indeslutter den ikke som det blot Afspeilede, men
som det, der kun har Væren ved det Væsen, der i sin Viden
omfatter den.

% p. 57: signature set right in heavy (bold) Fraktur, spelt „Bröchner“ with ö as printed; then a centred double rule (c. 3 cm) closes the article
\begin{flushright}\textbf{H. Bröchner.}\end{flushright}
\begin{center}\rule{3cm}{0.4pt}\\[-10pt]\rule{3cm}{0.4pt}\end{center}

\end{document}
